\chapter{Smith's Phase Transition and the Density of Attractors}
\label{sec:phase}
% ===================================================================

\section{Life as a Fourth Geosphere}

Eric Smith proposes that life constitutes a \emph{fourth geosphere}
--- alongside the lithosphere, hydrosphere, and atmosphere ---
that emerged from prebiotic chemistry through a phase transition.
The key mechanism is \emph{autocatalysis}: certain chemical reaction
networks catalyze their own production.
Once such a network is initiated, it becomes self-reinforcing,
crossing a thermodynamic threshold beyond which it is stable against
fluctuation.
The emergence of life is the crossing of this threshold.

\section{The Density of Recursive Attractors}

The connection to the present framework runs through the logical
topology of Part~II and the fixed-point structure of
Section~\ref{sec:replication}.

\begin{conjecture}[Density of Replicating Terms]
\label{conj:density}
The set of \emph{replicating terms} in the Rho calculus,
\[
  \mathrm{Rep} \;=\;
  \bigl\{ [P] \in \terms(\Rho)/{\bisim} \;\mid\;
    P \bisim P \mid P' \;\text{ for some }\; P' \bisim P \bigr\},
\]
is \emph{dense} in the logical topology $\mathcal{O}_{\HML}$:
every non-empty open set $\llb\phi\rrb \subseteq \terms(\Rho)/{\bisim}$
contains a replicating term.
\end{conjecture}

\begin{remark}[The Argument for Density]
Given any process $P$, one can construct a process $\hat{P}$
bisimilar to $P$ on all observable behaviors but additionally
capable of replication.
The construction uses the concurrent Y combinator to add
a replication sub-process running in parallel: 
$\hat{P} = P \mid \mathbf{Y}(\mathrm{copy})$
where $\mathrm{copy}$ is a process that duplicates $P$ onto a
fresh channel.
The added component runs in parallel and does not interfere with
$P$'s observable interactions, so $\hat{P}$ satisfies all the
same HML formulae as $P$ --- except formulae that specifically
probe the replication channel.
Since every open set $\llb\phi\rrb$ is defined by a single formula
$\phi$, and since $\phi$ generically does not probe the replication
channel, $\hat{P} \in \llb\phi\rrb$.
Making this argument rigorous requires showing that the added
replication machinery can always be hidden from the formula $\phi$
--- a careful but plausible argument deferred to future work.
\end{remark}

\section{Phase Transition as Basin Crossing}

\begin{remark}[The Prebiotic Wandering]
Consider a prebiotic population of processes in the Rho calculus,
wandering through the bisimulation quotient $\terms(\Rho)/{\bisim}$
under selective pressure.
Interactions that succeed --- those with high path amplitude
$W(\gamma)$ and affordable cost $\vect{c} \leq \vect{A}$ ---
are reinforced.
Interactions that fail deplete the vectorial account.
The population drifts under this selection pressure without
direction or goal.
\end{remark}

\begin{remark}[Inevitable Entry into the Basin]
By the density of $\mathrm{Rep}$ (Conjecture~\ref{conj:density}),
every open neighborhood in the logical topology contains a
replicating term.
The wandering population therefore inevitably enters a neighborhood
of $\mathrm{Rep}$.
This is not a low-probability event requiring fine-tuning; it is
structurally guaranteed by density.
The question is not \emph{whether} the population encounters a
replicator but \emph{when}.
\end{remark}

\begin{remark}[The Phase Transition]
Once the population enters the basin of attraction of a replicating
fixed point, the dynamics change qualitatively.
A replicating process produces copies of itself, increasing
$\mathrm{CN}(P, \mathcal{N})$.
More copies mean more interactions, more reinforcement, higher
path amplitudes for the replication pathway, lower effective cost
per replication event.
The population converges onto the replicating fixed point.

What emerges from the funnel is a population with a qualitatively
new feature: \emph{hereditary replication with variation} ---
the precondition for Darwinian evolution.
This is Smith's phase transition, made precise in the GSLT
framework.

The new geosphere, in our terms, is the \emph{compact coherence
cluster} (Part~II, Section~6) centered on the replicating fixed
point: the maximal compact subpopulation within which the
replicating process and its variants can reach consensus.
The phase transition is the crossing of the basin boundary in
the logical topology --- the moment at which the subpopulation
containing the replicator becomes compact, and consensus (here,
consensus on the replication strategy) becomes achievable.
\end{remark}

\begin{remark}[The Joint Rise of Assembly Index and Copy Number]
After the phase transition, both $\mathrm{AI}$ and $\mathrm{CN}$
increase together.
Selection amplifies processes that replicate more efficiently,
which typically requires more sophisticated catalytic machinery
--- higher $\mathrm{AI}$.
And more efficient replication produces more copies --- higher
$\mathrm{CN}$.
The joint biosignature of Assembly Theory is therefore not merely
an empirical observation but a \emph{theorem about the
post-transition dynamics} near a replicating fixed point, given
density of attractors and account-governed selection pressure.
This is a key prediction of the framework: wherever a phase
transition of this kind occurs, the joint $\mathrm{AI}$-$\mathrm{CN}$
signature will be observed.
\end{remark}

% ===================================================================
